% ID: SEL-D0006-OA-E04
% ESTADO: cotejada. Segunda pasada visual completa y borrador compilado revisados.
% Fuente: Selectividad/sel_2019_matematicas_aplicadas/Reserva B/Exámenes Matemáticas Aplicadas CCSS - A y B.pdf
% SHA-256: f9236ec66af0dbd3fbd0231fdcaa6d922425c04d48178ea0737073e76e8eb7d5
% Página: 1; opción A; ejercicio 4. Se conserva el dato 120 del original.

La producción en kilogramos por árbol de aguacates de una comarca sigue una
distribución Normal de desviación típica 4 y media desconocida.
\begin{enumerate}[label=\alph*)]
  \item \textbf{(1 punto)} Obtenga el tamaño muestral mínimo necesario para
  estimar la media poblacional con un error de estimación inferior a 2.1 kg
  y una confianza del 97\,\%.
  \item \textbf{(1.5 puntos)} Se toma una muestra aleatoria de 9 árboles,
  cuyas producciones en kilogramos han sido:
  \[
    15\quad 120\quad 50\quad 40\quad 5\quad 46\quad 52\quad 48\quad 10
  \]
  Obtenga el intervalo de confianza al 97\,\% para estimar la producción
  media de aguacates por árbol y calcule el error máximo de estimación.
\end{enumerate}
